\begin{proof}
This proof is a verbatim reproduction of that in \cite{AssadiCK19}. For any $i \in [m]$, define $\bA^{<i} := \{A_{x_1},\ldots,A_{x_{i-1}}\}$, $\bB^{>i} := \{B_{x_{i+1}},\ldots,B_{x_{m}}\}$. Recall that 
the internal information cost of $\prot_\SI$ is $\IC{\prot_\SI}{\cD_\SI} := \mi{\rA}{\rProtSI \mid \rB} + \mi{\rB}{\rProtSI \mid \rA}$. In the following, we focus on bounding the first term. The second term can be bounded exactly
the same by symmetry.  

As $(\rI,\bA^{<\rI},\bB^{>\rI},\bC,\bD)$ is sampled via public randomness in $\prot_\SI$, by Proposition~\ref{prop:public-random},   
\begin{align*}
	\mi{\rA}{\rProtSI \mid \rB} &= \mi{\rA}{\rProtSI \mid \rB, \rI,\bA^{<\rI},\bB^{>\rI},\bC,\bD} \leq \mi{\rA}{\rProtSI \mid \rB, \rI,\bA^{<\rI},\bB^{>\rI}}. 
\end{align*}
The inequality is by Proposition~\ref{prop:info-decrease} as we now show $\rA \perp (\bC,\bD) \mid \rProtSI,\rB,\rI,\bA^{<\rI},\bB^{>\rI}$ (and hence conditioning on $(\bC,\bD)$ can only decrease the mutual information). 
This is because $\rA \perp (\bC,\bD) \mid \rB,\rI,\bA^{<\rI},\bB^{>\rI}$ by Observation~\ref{obs:distHPC} and $\rProtSI$ is transcript of a deterministic protocol plus $z_1,\ldots,z_j$ obtained deterministically and hence we can apply
Proposition~\ref{prop:hpc-rectangle}.  

Define a random variable $\rTheta \in \{0,1\}$ where $\rTheta = 1$ iff in Line~(\ref{line:AB}) of protocol $\prot_\SI$, we terminate the protocol. In other words $\rTheta = 1$ iff $x_i \in Z^{<j}$. Since $\rA \perp \rTheta \mid \rB, \rI,\bA^{<\rI},\bB^{>\rI}$, 
further conditioning on $\rTheta$ can only increase the mutual information term above by Proposition~\ref{prop:info-increase}, hence, 
\begin{align}
	\mi{\rA}{\rProtSI \mid \rB} &\leq \mi{\rA}{\rProtSI \mid \rB, \rI,\bA^{<\rI},\bB^{>\rI},\rTheta} \notag \\
	&= \frac{m-j}{m} \cdot \mi{\rA}{\rProtSI \mid \rB, \rI,\bA^{<\rI},\bB^{>\rI},\rTheta=0} + \frac{j}{m} \cdot \mi{\rA}{\rProtSI \mid \rB, \rI,\bA^{<\rI},\bB^{>\rI},\rTheta=1} \notag \\
	&\leq \frac{m-j}{m} \cdot \mi{\rA}{\rProtSI \mid \rB, \rI,\bA^{<\rI},\bB^{>\rI},\rTheta=0}, \label{eq:mi-ra}
\end{align}
since conditioned on $\rTheta=1$, the protocol $\rProtSI$ is simple some prefix of $\rZ^{<j}$ and is hence independent of the input $(\rA,\rB)$ and carries no information about $\rA$ (see \itfacts{info-zero}). 
We now further bound the RHS of Eq~(\ref{eq:mi-ra}). When $\rTheta = 0$, $\rProtSI = (\rZ^{<j} ,\rProt_1,\ldots,\rProt_j) = (\rE^{<j},\rProt_j)$. Hence, we can write, 
  \begin{align*}
	 \mi{\rA}{\rProtSI \mid \rB, \rI,\bA^{<\rI},\bB^{>\rI},\rTheta=0} &\leq \mi{\rA}{\rE_j,\rProt_j \mid \rB, \rI,\bA^{<\rI},\bB^{>\rI},\rTheta=0} \\
	&= \mi{\rA}{\rZ^{<j} \mid \rB, \rI,\bA^{<\rI},\bB^{>\rI},\rTheta=0} \\
	&\hspace{25pt} + \mi{\rA}{\rProt^{<j},\rProt_j \mid \rZ^{<j},\rB, \rI,\bA^{<\rI},\bB^{>\rI},\rTheta=0} \tag{by chain rule in \itfacts{chain-rule} and since $\rE_j = (\rProt^{<j},\rZ^{<j})$} \\
	&\leq \mi{\rA}{\rProt \mid \rZ^{<j},\rB, \rI,\bA^{<\rI},\bB^{>\rI},\rTheta=0},
\end{align*}
as $\rA \perp \rZ^{<j} \mid \rTheta=0$ (and other variables) and hence the first term is zero, and in the second term $\rProt$ contains $\rProt^{<j},\rProt_j$ (plus potentially other terms) and so having $\rProt$ in instead can only increase
the information. By further expanding the conditional information term above,
\begin{align*}
	 &\mi{\rA}{\rProtSI \mid \rB, \rI,\bA^{<\rI},\bB^{>\rI},\rTheta=0} \\
	 &\hspace{50pt}\leq \E_{(Z^{<j},i) \mid \rTheta=0} \bracket{\mi{\rA}{\rProt \mid \rB, \bA^{<i},\bB^{>i},\rI=i,\rZ^{<j}=Z^{<j},\rTheta=0}} \\
	&\hspace{50pt}= \E_{Z^{<j} \mid \rTheta=0} \bracket{\sum_{\substack{i=1 \\~i \notin Z^{<j}}}^{m} \frac{1}{m-j}  \mi{\rA_{x_i}}{\rProt \mid \rB_{x_i}, \bA^{<i},\bB^{>i},\rI=i,\rZ^{<j}=Z^{<j},\rTheta=0}} \tag{conditioned on $\rTheta=0$, $i$ is chosen
	uniformly at random from $Z^{<j}$; also $(\rA,\rB) = (\rA_{x_i},\rB_{x_i})$} \\
	&\hspace{50pt}= \E_{Z^{<j} \mid \rTheta=0} \bracket{\sum_{i \notin Z^{<j}} \frac{1}{m-j} \cdot \mi{\rA_{x_i}}{\rProt \mid \rB_{x_i}, \bA^{<i},\bB^{>i},\rZ^{<j}=Z^{<j},\rTheta=0}}
	\tag{we dropped the conditioning on $\rI=i$ as all remaining variables are independent of this event}  \\
	&\hspace{50pt}= \E_{Z^{<j} \mid \rTheta=0} \bracket{\sum_{i \notin Z^{<j}} \frac{1}{m-j} \cdot \mi{\rA_{x_i}}{\rProt \mid \bA^{<i},\bB,\rZ^{<j}=Z^{<j},\rTheta=0}}
	\tag{as $\rA_{x_i} \perp \bB^{<i} \mid \rB_{x_i},\bA^{<i}$ by Observation~\ref{obs:distHPC} and hence we can apply Proposition~\ref{prop:info-increase}} \\
	&\hspace{50pt}\leq \E_{Z^{<j} \mid \rTheta=0} \bracket{\sum_{i=1}^{m} \frac{1}{m-j} \cdot \mi{\rA_{x_i}}{\rProt \mid \bA^{<i},\bB,\rZ^{<j}=Z^{<j},\rTheta=0}} \\ 
	\tag{mutual information is non-negative by \itfacts{info-zero} and so we can add the terms in $Z^{<j}$  as well} \\
	&\hspace{50pt}= \E_{Z^{<j} \mid \rTheta=0} \bracket{\sum_{i=1}^{m} \frac{1}{m-j} \cdot \mi{\rA_{x_i}}{\rProt \mid \bA^{<i},\bB,\rZ^{<j}=Z^{<j},\rTheta=0}} \\
	&\hspace{50pt}= \frac{1}{m-j} \cdot \E_{Z^{<j} \mid \rTheta=0} \bracket{\mi{\bA}{\rProt \mid \bB,\rZ^{<j}=Z^{<j},\rTheta=0}} \tag{by chain rule in \itfacts{chain-rule}} \\
	&\hspace{50pt}= \frac{1}{m-j} \cdot \mi{\bA}{\rProt \mid \bB,\rZ^{<j},\rTheta=0} \\
	&\hspace{50pt}\leq \frac{1}{m-j} \cdot \paren{\mi{\bA}{\rProt \mid \bB,\rTheta=0} + \en{\rZ^{<j}}} \\
	&\hspace{50pt}= \frac{1}{m-j} \cdot \paren{\mi{\bA}{\rProt \mid \bB} + \en{\rZ^{<j}}} \tag{transcript of the protocol $\prot_\bhpc$ (namely $\rProt$) on input $(\bA,\bB)$ is independent of $\rTheta$} \\
	&\hspace{50pt}\leq \frac{1}{m-j} \cdot \paren{\en{\rProt} + \en{\rZ^{<j}}} \leq \frac{\CC{\prot_\bhpc}{} }{m-j}   + \frac{j \cdot \log{m}}{m-j}. \tag{by sub-additivity of entropy (\itfacts{sub-additivity}) and \itfacts{uniform}}
\end{align*}

By plugging in this bound in Eq~(\ref{eq:mi-ra}), we have that,
\begin{align*}
	\mi{\rA}{\rProtSI \mid \rB} &\leq \frac{m-j}{m} \cdot \paren{\frac{\CC{\prot_\bhpc}{} }{m-j}   + \frac{j \cdot \log{m}}{m-j}} = \frac{\CC{\prot_\bhpc}{} }{m}   + \frac{j \cdot \log{m}}{m}. 
\end{align*}
By symmetry, we can also prove the same bound on $\mi{\rB}{\rProtSI \mid \rA}$. As such, we have,
\begin{align*}
	\mi{\rA}{\rProtSI \mid \rB} + \mi{\rB}{\rProtSI \mid \rA} \leq 2 \cdot \paren{\frac{\CC{\prot_\bhpc}{} }{m}   + \frac{j \cdot \log{m}}{m}}.
\end{align*}

We shall note that strictly speaking the factor $2$ above is not needed (similar to the proof of Proposition~\ref{prop:cc-ic}) but as this factor is anyway suppressed through O-notation later in the proof, the above bound
suffices for our purpose. %\qed\; %\Cref{clm:direct-sum}

\end{proof}
